%!TEX root = main.tex
\setcounter{chapter}{4}
\setcounter{section}{3}
\setcounter{theorem}{0}
\setcounter{equation}{0}

\lectureheader{161}{Calculus I}{First derivative test}{\textit{Thomas' Calculus} \thesection}

\begin{definition}
A function is said to be \textbf{monotonic} on an interval if it is increasing on the interval or decreasing on the interval.
\end{definition}

\begin{theorem}
Suppose that $f	$ is continuous on $[a,b]$ and differentiable on $(a,b)$.
\begin{itemize}
\item If $f'(x)>0$ for each $x\in (a,b)$, then $f$ is increasing on $[a,b]$.
\item If $f'(x)<0$ for each $x\in (a,b)$, then $f$ is decreasing on $[a,b]$.
\end{itemize}
\end{theorem}
\begin{proof}\,
\vspace{5.5in}

\end{proof}

\newpage

\begin{corollary}[First derivative test]
Suppose that $c$ is a critical point of a continuous function $f$, and that $f$ is differentiable on a open interval containing $c$ except possibly at $c$.
\begin{enumerate}
\item If $f'(x)$ changes from negative to positive at $x=c$, then $f$ has a local minimum at $x=c$.
\item If $f'(x)$ changes from positive to negative at $x=c$, then $f$ has a local maximum at $x=c$.
\item If $f'(x)$ does not change sign at $x=c$, then $f$ does not have a local extremum at $x=c$.
\end{enumerate}
\end{corollary}

\begin{example}
Identify the open intervals on which $f(x)=x^3-12x-5$ is increasing and on which $f$ is decreasing.
Also identify the function's local and global extrema.
\end{example}

\newpage

\begin{example}
Identify the open intervals on which $f(x)=x^{1/3}(x-4)$ is increasing and on which $f$ is decreasing.
Also identify the function's local and global extrema.
\end{example}

\newpage

\begin{example}
Consider $f(x)=\sin^2 x -\sin x-1$ on the interval $[0,2\pi]$.
Identify the open intervals on which $f$ is increasing and on which $f$ is decreasing.
Also identify the function's local and global extrema.
\end{example}
